\documentclass[a4paper, 12pt]{report}
% Packages
\usepackage{amsmath}
\usepackage{amssymb}
\usepackage[margin=1in]{geometry}
\usepackage{hyperref}
\usepackage{graphicx}
\usepackage{subcaption}
\usepackage{float}
\usepackage{array}
\usepackage{booktabs}

%Indentation and spacing
\setlength{\parindent}{0pt}

\begin{document}

%%%%%% Page - 1 %%%%%%
\section*{Multi-Frequency method}

\begin{flushleft}
This section is organized as follows: \\
\hspace{2cm}1. Identification of aerodynamic forces \\
\hspace{2cm}2. Calculation of Flutter Derivatives from aerodynamic forces \\
\hspace{2cm}3. Calculation of Critical Wind Speed from Flutter Derivatives \\
\hspace{2cm}4. Multi-Frequency method \\
\end{flushleft}

\subsection*{Identification of aerodynamic forces:}

Obtaining Flutter Derivatives with Forced-Motion CFD simulations:

\begin{itemize}
    \setlength{\itemsep}{10pt}
    \item[$\bullet$] 2-D cross-section of the bridge is modelled in Virtual Wind Tunnel with appropriate BC

    \item[$\bullet$] The bridge cross-section is oscillated sinusoidally in the vertical direction (Heave motion)
    \begin{align}
        h(t)        &= h_0 \sin(\omega t) \\
        \dot{h}(t)  &= \omega h_0 \cos(\omega t)
    \end{align}
    where, $h_0$ - Maximum amplitude of the applied heave motion \\
    \phantom{where,} $\omega$ - Frequency (real number) of the applied heave motion

    \item[$\bullet$] Since pitch motion is not applied, $\alpha$ and its derivative $= 0$. Hence, the lift and moment force can be written as
    \begin{align}
        L_{ae} &= \frac{1}{2}\rho U^2 B \left[KH_1^* \frac{\dot{h}}{U} + K^2H_4^* \frac{h}{B}\right] \quad \text{(From Scanlan Formulation)} \\
        M_{ae} &= \frac{1}{2}\rho U^2 B^2 \left[KA_1^* \frac{\dot{h}}{U} + K^2A_4^* \frac{h}{B}\right] \quad \text{(From Scanlan Formulation)}
    \end{align}

    \item[$\bullet$] Then, the aerodynamic lift force and moment acting on the cross-section are measured from CFD. As per Scanlan's assumption, under small amplitude (for $h_0$), aerodynamic force will be
    \begin{align}
        L_{ae} &= B_0 \sin(\omega t + \phi) \quad \text{(From CFD)} \\
        L_{ae} &= B_1 \cos(\omega t) + B_2 \sin(\omega t)
    \end{align}
    \begin{align}
        M_{ae} &= C_1 \sin(\omega t + \Phi) \quad \text{(From CFD)} \\
        M_{ae} &= C_1 \cos(\omega t) + C_2 \sin(\omega t)
    \end{align}
\end{itemize}

\vspace{0.5cm}

\subsection*{Calculation of Flutter Derivatives:}

\begin{itemize}
    \setlength{\itemsep}{10pt}
    \item[$\bullet$] On equating the lift force from Scanlan's formulation and CFD
    \begin{equation}
        \frac{1}{2}\rho U^2 B \left[KH_1^* + K^2H_2^*\right] = H_1 \cos(\omega t) + H_2 \sin(\omega t)
    \end{equation}
    \begin{equation}
        \frac{1}{2}\rho U^2 B^2 \left[KA_1^* + K^2A_2^*\right] = C_1 \cos(\omega t) + C_2 \sin(\omega t)
    \end{equation}
    Hence,
    \begin{align}
        H_1^* &= \frac{2B_1}{\rho \omega^2 B^2 h_0} & H_2^* &= \frac{2B_2}{\rho \omega^2 B^2 h_0} \\
        A_1^* &= \frac{2C_1}{\rho \omega^2 B^3 h_0} & A_2^* &= \frac{2C_2}{\rho \omega^2 B^3 h_0}
    \end{align}

    \item[$\bullet$] The same procedure is done with the Pitch motion to obtain $H_3^*$, $H_4^*$, $A_3^*$ and $A_4^*$

    \item[$\bullet$] Hence, 2 simulations are required to obtain 1 set of flutter derivatives
\end{itemize}

\vspace{0.5cm}

\subsection*{Calculation of Critical Wind Speed:}

Equations of motion for structural components
\begin{align}
    m(\ddot{h} + 2\omega_{h}\zeta_h \dot{h} + \omega_h^2 h)                   &= L_{ae} \\
    I_\alpha(\ddot{\alpha} + 2\omega_{\alpha}\zeta_{\alpha} \dot{\alpha} + \omega_{\alpha}^2 \alpha) &= M_{ae}
\end{align}

Aerodynamic forces
\begin{align}
    L_{ae} &= \frac{1}{2}\rho U^2 B \left[KH_1^* \frac{\dot{h}}{U} + KH_2^* \frac{B\dot{\alpha}}{U} + K^2H_3^* \alpha + K^2H_4^* \frac{h}{B}\right] \\
    M_{ae} &= \frac{1}{2}\rho U^2 B^2 \left[KA_1^* \frac{\dot{h}}{U} + KA_2^* \frac{B\dot{\alpha}}{U} + K^2A_3^* \alpha + K^2A_4^* \frac{h}{B}\right]
\end{align}

On substituting aerodynamic forces in the equations of motion
\begin{align}
    m(\ddot{h} + 2\omega_h\zeta_h \dot{h} + \omega_h^2 h)
    - \frac{1}{2}\rho U^2 B \left[KH_1^* \frac{\dot{h}}{U} + KH_2^* \frac{B\dot{\alpha}}{U}
    + K^2H_3^* \alpha + K^2H_4^* \frac{h}{B}\right] &= 0 \\
    I_\alpha(\ddot{\alpha} + 2\omega_{\alpha}\zeta_{\alpha} \dot{\alpha} + \omega_{\alpha}^2 \alpha)
    - \frac{1}{2}\rho U^2 B^2 \left[KA_1^* \frac{\dot{h}}{U} + KA_2^* \frac{B\dot{\alpha}}{U}
    + K^2A_3^* \alpha + K^2A_4^* \frac{h}{B}\right] &= 0
\end{align}

Writing it in matrix format in time domain,
\begin{equation}
\begin{split}
    &\begin{bmatrix} m & 0 \\ 0 & I \end{bmatrix}
    \begin{Bmatrix} \ddot{h} \\ \ddot{\alpha} \end{Bmatrix}
    + \begin{bmatrix} C_h & 0 \\ 0 & C_{\alpha} \end{bmatrix}
    \begin{Bmatrix} \dot{h} \\ \dot{\alpha} \end{Bmatrix}
    + \begin{bmatrix} K_h & 0 \\ 0 & K_{\alpha} \end{bmatrix}
    \begin{Bmatrix} h \\ \alpha \end{Bmatrix} \\
    &\quad - \frac{1}{2}\rho K^2
    \left(
    \begin{bmatrix} H_1^* & BH_2^* \\ A_1^* & BA_2^* \end{bmatrix}
    \begin{Bmatrix} \dot{h} \\ \dot{\alpha} \end{Bmatrix}
    + \begin{bmatrix} H_4^* & H_3^* \\ A_4^* & BA_3^* \end{bmatrix}
    \begin{Bmatrix} h \\ \alpha \end{Bmatrix}
    \right)
    = \begin{Bmatrix} 0 \\ 0 \end{Bmatrix}
\end{split}
\end{equation}

Converting to the frequency domain with $h = e^{i\omega t}$ and $\alpha = e^{i\omega t}$
\begin{multline}
    \left(
    -\omega^2 \begin{bmatrix} m & 0 \\ 0 & I \end{bmatrix}
    + \frac{1}{2}\rho K^2
    \left(
    \begin{bmatrix} H_1^* & BH_2^* \\ A_1^* & BA_2^* \end{bmatrix}
    - \begin{bmatrix} H_4^* & BH_3^* \\ A_4^* & BA_3^* \end{bmatrix}
    \right)
    + i\omega \begin{bmatrix} C_h & 0 \\ 0 & C_\alpha \end{bmatrix}
    + \begin{bmatrix} K_h & 0 \\ 0 & K_\alpha \end{bmatrix}
    \right)
    \begin{Bmatrix} h \\ \alpha \end{Bmatrix}
    = \begin{Bmatrix} 0 \\ 0 \end{Bmatrix}
\end{multline}

Where $\omega = \omega_r + i\omega_i$, the real part of $\omega$ represents the oscillating frequency and the imaginary part of $\omega$ represents damping.

This is in the form
\begin{equation}
    -\omega^2 M + i\omega C + K = 0
\end{equation}
where,
\begin{align}
    M &= \begin{bmatrix} m & 0 \\ 0 & I \end{bmatrix}
         + \frac{1}{2}\rho K^2
         \left(
         \begin{bmatrix} H_1^* & BH_2^* \\ A_1^* & BA_2^* \end{bmatrix}
         - \begin{bmatrix} H_4^* & BH_3^* \\ A_4^* & BA_3^* \end{bmatrix}
         \right) \\
    C &= \begin{bmatrix} C_h & 0 \\ 0 & C_\alpha \end{bmatrix} \\
    K &= \begin{bmatrix} K_h & 0 \\ 0 & K_\alpha \end{bmatrix}
\end{align}

Since damping is present, our equation takes the form of a complex quadratic eigenvalue problem. Hence, it is linearized and solved for $\omega$.

Linearization,
\begin{equation}
    \begin{bmatrix} 0 & K \\ K & iC \end{bmatrix}
    - \omega \begin{bmatrix} K & 0 \\ 0 & -M \end{bmatrix}
    \begin{pmatrix} \frac{X}{\omega X} \end{pmatrix}
    = \begin{bmatrix} 0 \\ 0 \end{bmatrix}
\end{equation}
\begin{equation}
    A - \lambda B = 0
\end{equation}

Depending on the imaginary part of $\omega$, the system can be stable or unstable
\begin{itemize}
    \setlength{\itemsep}{10pt}
    \item[$\bullet$] If it is positive: $\omega = \omega_r + i\omega_i$, \quad $h = h_0 e^{i(\omega_r + i\omega_i)t} = h_0 (e^{i\omega_r t} \cdot e^{-\omega_i t})$ --- Stable
    \item[$\bullet$] If it is negative: $\omega = \omega_r + i\omega_i$, \quad $h = h_0 e^{i(\omega_r + i\omega_i)t} = h_0 (e^{i\omega_r t} \cdot e^{+\omega_i t})$ --- Unstable
\end{itemize}

A Python tool is created to solve this linearized eigenvalue problem and validated with G. Diana benchmarks. To find the critical wind speed with this approach is an iterative process. First, a reduced velocity is selected. For that reduced wind speed, aerodynamic forces are taken from the flutter derivative functions and solved for $\omega$ with the tool created for solving critical wind. If it is stable, then the next reduced wind speeds are progressively selected until the system reaches its reduced critical wind speed (unstable system). Based on the reduced critical wind speed and real part of $\omega$, the actual critical wind speed is obtained.


\subsection*{Multi-Frequency method:}

For the multi-frequency method, instead of applying a single frequency for the forced motion oscillation in heave and pitch, 
multiple frequencies are applied simultaneously. Assuming that aerodynamic forces are linear, the aerodynamic forces will be 
the superposition of the aerodynamic forces obtained from each frequency. Hence, with this approach, multiple sets of flutter 
derivatives can be obtained from a single simulation. Hence the section "Identification of aerodynamic forces" alone is different
on comparing the single frequency and multi-frequency method. The rest of the procedure for calculating flutter derivatives 
and critical wind speed is the same for both approaches. Figure \ref{fig:Single_Frequency_Approach} and 
\ref{fig:Multi-Frequency_Method} show the difference in the procedure for both approaches.

\begin{figure}[htbp]
    \centering
    \includegraphics[width=1\textwidth]{Chap_1/Figures/Single_Frequency.png}
    \caption{Single Frequency Approach}
    \label{fig:Single_Frequency_Approach}
\end{figure}



\begin{figure}[htbp]
    \centering
    \includegraphics[angle=0, width=1.0\textwidth]{Chap_1/Figures/MultiFrequency_Explanation.png}
    \caption{Multi-Frequency method}
    \label{fig:Multi-Frequency_Method}
\end{figure}

\end{document}